A crescente demanda por processamento em aplicações veiculares impulsionou a adoção de \gls{VEC}. Nesse paradigma, veículos com capacidade computacional restrita podem superar suas limitações locais através do descarregamento (\textit{offloading}) de tarefas para servidores de borda e veículos vizinhos. Contudo, otimizar as decisões de processamento sob restrições rigorosas de prazo e consumo energético é um desafio central devido à não estacionariedade desses ambientes. A mudança contínua no número de nós candidatos e a variação na carga de trabalho dificultam o uso direto de modelos tradicionais de aprendizado de máquina, que frequentemente exigem arquiteturas de entrada de dimensão fixa e retreinamento constante frente a essas variações. Para superar essas limitações, este trabalho propõe e avalia uma arquitetura híbrida de decisão baseada na integração de métodos multicritério com \gls{DRL}. O componente multicritério pondera de forma contínua aspectos mais críticos da tomada de decisão, como o tempo estimado de conclusão, o custo energético e a distância, fornecendo uma representação compacta do estado observado pelo agente enquanto o \gls{DRL} aprende qual a melhor política de descarregamento a longo prazo. Essa solução híbrida garante invariância à escala e estabilidade de aprendizado, permitindo que o agente inteligente opere sem colapso de política independentemente da densidade veicular. Avaliações em ambientes de simulação demonstram que a arquitetura proposta alcança uma taxa de cumprimento de prazos operacionais de 45,21\%, superando expressivamente múltiplas heurísticas de referência e outras estratégias fundamentadas em \gls{DRL}. Tais resultados atestam a eficiência sistêmica do modelo para cenários complexos. Como extensão desta proposta, será incorporada lógica de nebulosa (\textit{Fuzzy}) aos critérios de decisão, com o objetivo de aprimorar a resiliência da solução frente a informações imprecisas.

\palavraschave{Descarregamento de Tarefas. Redes Veiculares. Computação de Borda Veicular. Comunicação Veicular. Aprendizado por Reforço Profundo. Fuzzy-Topsis. Decisão Multicritério.}
